\documentclass[11pt,a4paper]{article}
\usepackage[utf8]{inputenc}
\usepackage[UTF8]{ctex} % 支持中文
\usepackage{geometry}
\geometry{left=2.5cm,right=2.5cm,top=2.5cm,bottom=2.5cm}
\usepackage{graphicx}
\usepackage{listings}
\usepackage{xcolor}
\usepackage{hyperref}
\usepackage{amsmath}
\usepackage{tcolorbox}

% 代码样式设置
\lstset{
    basicstyle=\ttfamily\small,
    keywordstyle=\color{blue},
    commentstyle=\color{gray},
    stringstyle=\color{orange},
    breaklines=true,
    frame=single,
    backgroundcolor=\color{lightgray!10}
}

\title{\textbf{ACE\_Agent 聚类智能体项目：从零开始的学习笔记}}
\author{Gemini CLI 助教}
\date{2026年4月13日}

\begin{document}

\maketitle

\begin{abstract}
本文档旨在为 Python 和聚类知识薄弱的学习者提供 ACE\_Agent 项目的深度解析。我们将从系统架构、工作流逻辑、核心代码实现以及聚类算法原理四个维度，带你全面掌握这个自动化聚类集成系统。
\end{abstract}

\tableofcontents

\newpage

\section{第一章：Python 基础扫盲}

在深入代码之前，必须掌握以下三个核心概念。

\subsection{1.1 包 (Package) 与 \texttt{\_\_init\_\_.py}}
在 Python 中，文件夹并不自动等于“模块”。
\begin{itemize}
    \item \textbf{概念}：包含 \texttt{\_\_init\_\_.py} 的文件夹被称为“包”。
    \item \textbf{作用}：它告诉 Python 该目录可以被 \texttt{import}。例如，\texttt{from agent\_core import router}。
\end{itemize}

\subsection{1.2 类 (Class) 与 构造器 (\texttt{\_\_init\_\_})}
类是建筑图纸，对象是真实的房子。
\begin{itemize}
    \item \textbf{构造器}：即 \texttt{def \_\_init\_\_(self):} 函数。当你创建对象（如 \texttt{supervisor = ACESupervisor()}）时，该函数会自动运行。
    \item \textbf{self}：代表对象自己，用于在不同函数之间传递“器官”（变量）。
\end{itemize}

\subsection{1.3 装饰器 (Decorator, \texttt{@...})}
\begin{itemize}
    \item \textbf{形象理解}：给函数穿上一件“功能外套”。
    \item \textbf{示例}：\texttt{@dataclass} 自动为类生成数据存储功能，无需手动写繁琐的初始化代码。
\end{itemize}

\section{第二章：ACE\_Agent 整体架构}

ACE\_Agent 模拟了一个数据科学团队的工作模式：

\begin{tcolorbox}[title=团队角色分配]
\begin{enumerate}
    \item \textbf{总监 (Supervisor)}：负责整体流程调度、记忆管理。
    \item \textbf{调度员 (Router)}：通过统计学指标对数据进行“体检”，决定派谁上场。
    \item \textbf{专家团 (Experts)}：质心专家、拓扑专家、维度专家等。
    \item \textbf{实验室 (Sandbox)}：安全地运行专家动态生成的代码。
    \item \textbf{工厂 (Data Factory)}：生产各种形状的数据集（笑脸、月牙等）。
    \item \textbf{秘书 (LaTeX Generator)}：将结果排版成精美的 PDF 报告。
\end{enumerate}
\end{tcolorbox}

\section{第三章：核心工作流详解}

\subsection{3.1 意图识别 (Intent Recognition)}
系统会首先判断用户是在问“新问题”还是在针对上一个结果“追问”。如果是追问，它会调用 LLM 的推理能力，而不会重新运行耗时的算法。

\subsection{3.2 数据画像 (Data Profiling)}
调度员会计算以下指标：
\begin{itemize}
    \item \textbf{偏度 (Skewness)}：判断数据分布是否对称。
    \item \textbf{非凸性 (Non-convexity)}：判断数据是否具有奇怪的形状（如月牙）。
    \item \textbf{相关性 (Correlation)}：判断特征之间是否冗余。
\end{itemize}

\subsection{3.3 专家协作与并行化}
为了提高效率，总监使用了 \textbf{线程池 (ThreadPoolExecutor)}。这意味着多个专家可以同时在各自的“办公室”里跑算法，互不干扰。

\section{第四章：数据工厂与 LaTeX 报告}

\subsection{4.1 数据生成原理}
项目使用了 \texttt{scikit-learn} 提供的合成数据生成器：
\begin{lstlisting}[language=Python]
# 生成两个月牙形状的数据
from sklearn.datasets import make_moons
X, y = make_moons(n_samples=400, noise=0.05)
\end{lstlisting}
这些数据被包装进 \texttt{DatasetBundle} 对象，流转到各个专家手中。

\subsection{4.2 报告自动生成}
秘书使用 LaTeX 模板，将 \texttt{AlgorithmRunResult} 中的字典数据转化为格式化的文本。
\begin{itemize}
    \item \textbf{表格}：汇总各算法的 AMI、轮廓系数。
    \item \textbf{图片}：插入各专家生成的聚类效果图。
    \item \textbf{结论}：由总监撰写执行摘要。
\end{itemize}

\section{第五章：聚类算法知识点}

\subsection{5.1 质心聚类 (K-Means)}
\textbf{原理}：不断寻找类中心，让点到中心的距离之和最小。\\
\textbf{局限}：只能处理圆形的簇，遇到月牙形会失败。

\subsection{5.2 密度聚类 (DBSCAN)}
\textbf{原理}：像传染病一样连通密度足够大的点。\\
\textbf{优点}：能识别任何形状，能自动发现噪声点。

\subsection{5.3 评价指标}
\begin{itemize}
    \item \textbf{轮廓系数 (Silhouette)}：无监督指标。1 代表完美，-1 代表分错。
    \item \textbf{AMI}：有监督指标。看算法结果与真实答案的相似度。
\end{itemize}

\section{结语}
ACE\_Agent 不仅仅是一个聚类工具，它展示了如何通过 \textbf{LLM + 规则引擎 + 自动化代码生成} 来构建一个智能数据科学助手。

\end{document}
